\chapter{Przegląd literatury}

\section{Markowskie procesy decyzyjne}

Teoria markowskich procesów decyzyjnych (MDP) stanowi fundamentalną ramę teoretyczną dla modelowania sekwencyjnych problemów decyzyjnych w warunkach niepewności.

\section{Dyskontowanie quasi-hiperboliczne}

Dyskontowanie quasi-hiperboliczne zostało wprowadzone w celu lepszego modelowania preferencji czasowych w sytuacjach, gdzie tradycyjne dyskontowanie wykładnicze nie oddaje rzeczywistych zachowań decydentów.

\subsection{Literatura podstawowa}

\cite{bertsekas2019} przedstawia kompleksowe wprowadzenie do uczenia ze wzmocnieniem i optymalnego sterowania.

\cite{jaskiewicz2021} wprowadzają markowskie procesy decyzyjne z dyskontowaniem quasi-hiperbolicznym, co stanowi bezpośrednią podstawę teoretyczną dla niniejszej pracy.

\cite{eshwar2024} rozwijają algorytmy uczenia ze wzmocnieniem dla dyskontowania quasi-hiperbolicznego, przedstawiając praktyczne metody implementacji.

\section{Niespójność czasowa}

Problem niespójności czasowej w procesach decyzyjnych jest kluczowym zagadnieniem, które motywuje zastosowanie dyskontowania quasi-hiperbolicznego zamiast tradycyjnego wykładniczego.